\section{Summary of Findings}
\label{sec:conclusion-summary}

This thesis asked how repository context changes review coverage, whether
structural context enables cross-file reasoning, and whether human raters
perceive that capability.

\subsection*{RQ1: Comparative coverage}

There is no single best review mode. Retrieval produces the largest gain on
the full rubric ($+1.03$, $p = 0.002$), while the structural
arm produces the largest gain on the nine structurally relevant criteria
($+0.69$, $p = 0.003$).
After correction for multiple comparisons, each retains the gain on the scale
on which it leads; the hybrid retains neither. The graph-arm estimate is
sensitive to regeneration and to the judge panel, so its magnitude should be
read with the uncertainty reported in Section~\ref{sec:threats}.

\subsection*{RQ2: Cross-file reasoning}

Structural context enables a capability that the diff-only reviewer lacks. On
28 defects whose consequences lie outside the edited file, the baseline
detects none and the graph arm detects 18 ($p < 0.0001$); retrieval
detects 5. The hybrid detects 21 but is not significantly better than the graph
arm alone. Adding inheritance edges raises detection to 26 of 28. These
results establish the capability on controlled
synthetic defects; they do not estimate its prevalence in the broader
population of pull requests.

\subsection*{RQ3: Human preference}

Twenty-seven raters preferred the graph-augmented review when asked which review names the affected code
(0.710, $p = 0.0005$). They showed no preference on
overall usefulness (0.512). After correction they preferred the baseline
review on code clarity. The human-study result therefore corroborates the
targeted affected-code benefit, not a general improvement in review quality.

\subsection*{Synthesis}

In the controlled experiment, structural context gives the language model a
capability unavailable from the diff alone: identifying a consequence that is
not in the diff. On the
35-pull-request benchmark its average gain is small but sharply localised:
96\% ($p = 0.003$) lands in the criteria the mechanism
predicts.

\section{Limitations}
\label{sec:conclusion-limitations}

This work has three boundaries. First, it models static relationships within
one repository, not dynamic or cross-repository dependencies. Second, the
evaluation uses selected pull requests, synthetic defects, and a small
Python-only human study. Third, it measures review coverage, controlled
defect detection, and preference rather than production defects or developer
productivity. Section~\ref{sec:threats} discusses the methodological threats
in detail.

\section{Future Work}
\label{sec:conclusion-future}

\paragraph{Issue-guided cross-repository graphs.}
The graphs in this thesis stop at the repository boundary and are seeded from
the changed files. Future work should connect related repositories through
APIs, schemas, package manifests, events, and test dependencies. Issue and
pull-request descriptions can identify relevant entities and intended
behaviour, then guide a demand-driven traversal of this larger graph. Textual
descriptions should rank and restrict the retrieved subgraph rather than create
structural edges by themselves. A controlled evaluation could compare a
single-repository graph, a complete cross-repository graph, and an issue-guided
cross-repository subgraph under the same context budget. Relevant outcomes
include detection of downstream breakage, false-positive rate, context size,
latency, and token cost.

\paragraph{Review-history memory.}
The current system treats each pull request independently. A temporal graph
could connect code entities to earlier issues, pull requests, review comments,
and the outcomes of those comments, including whether a finding was accepted,
fixed, or rejected. A reviewer could then retrieve relevant precedents when a
similar component or contract changes again. Evaluation should use a strict
temporal split, allowing access only to history that predates each test pull
request. Relevant outcomes include precision on later reviews, false-positive
reduction, developer acceptance, and adaptation when project conventions
change. Such a system would also need mechanisms for forgetting stale
preferences and avoiding the reinforcement of incorrect historical decisions.

\paragraph{Adaptive context selection.}
The pipeline supplies the same context block to every pull request, although
the graph arm's benefit is concentrated in changes with cross-file
consequences (Section~\ref{sec:discussion-practical}). A lightweight gate
could instead decide per pull request whether graph context is warranted.
The simplest gate is deterministic: compute the change's structural fan-out
from the code property graph before generation and supply the graph block
only when the changed functions have dependents outside the edited files. A
learned alternative would follow adaptive retrieval, where a small classifier
or the generator itself predicts whether retrieval will help and skips it
otherwise~\cite{jeong2024adaptiverag,asai2024selfrag,jiang2023flare};
Repoformer applies this to repository-level code completion and reports that
retrieval helps on only a minority of instances~\cite{wu2024repoformer}.
Pull-request change-impact analysis already derives per-file risk from
call-graph dependencies and history~\cite{gocmen2025changeimpact}, so the
gate signal is available at review time. Evaluation should follow the
routing literature: compare the gated pipeline with always-graph and
never-graph baselines on rubric score and detection, report tokens and
latency, and measure the gate's agreement with an oracle label marking the
pull requests on which the graph arm outperformed the
baseline~\cite{ong2025routellm,chen2024frugalgpt}. The risk is that a gate's
false negatives reinstate the blind spot the graph was introduced to close,
so recall on cross-file changes should be weighted above precision.

\paragraph{Correctness-oriented and cross-language evaluation.} A rubric that
asks whether a review's claims are \emph{grounded} (whether the component it
names is genuinely affected) would measure what matters. Cross-repository
contract tests and runtime traces could verify predicted breakage, while
developer acceptance or rejection of review findings would measure practical
utility. A follow-up human study should include pull requests from non-Python
repositories and language-appropriate graph builders to test whether the
affected-code preference and clarity trade-off generalise across languages.
Together, these evaluations would expose cases in which syntactically valid
but coarse-grained edges support irrelevant claims.

\paragraph{Human calibration of the judge panel.}
This thesis validates its LLM judges only against each other and never
against human judgement. The human study and Experiment~1 also cannot be
compared directly, since they differ in stimuli, graph construction and
criteria. A direct calibration check would close this gap at low cost. The
five-judge panel would score the twelve human-study reviews in the same
pairwise format the raters used, with the H1--H6 wording unchanged. For each
pull request and criterion, the panel's majority verdict would then be
compared with the raters' majority preference, and their agreement reported.
Each human criterion would also be mapped to its closest rubric criterion: H1
to F3 and F4, H2 to F2, H3 to T3, and H5 to R1. This tests whether the
effects point in the same direction in both instruments. Such a check would
show whether the panel reproduces human preferences, which is the
precondition for reading the 35-pull-request rubric results as evidence of
review quality. With only six pull requests, the agreement estimates would be
imprecise. The check should therefore be repeated on a larger stimulus set
built with the Joern pipeline.
